\section{Spirit in the World}

Within a traditional civilization, categorized as ``civilizations of space'' by \textbf{Julius Evola}, time seems to stand still. A citizen then would have had a much different experience of time than we do today. There was no notion of ``progress'', as Evola explains:

\begin{quotex} 
The civilizations of space , according to the formulation current today, would therefore be the ``stable'', ``static'' or ultraconservative civilizations. In reality, they are the civilizations whose material vestiges even seem destined to live longer than all the monuments or all the ideal creations of the modern world that appear powerless to endure more than a half century and within which the words ``progressive'' or ``dynamism'' mean only a compliance to contingency, to the movement of an incessant change, of a rapid ascent and a likewise rapid, vertiginous decline.
\end{quotex}


Thus, today a man must make a choice, a choice that was inconceivable to the original traditional man. That is because we today are aware of both worlds. Nevertheless, that decision opens up more possibilities and, since time is ``quicker'' today, more rapid results can be gained. First it is necessary to become detached from the verbal representation of the modern world in order to see it as just one worldview among many others. Then, through meditation and exercises, one can begin to re-experience the world in a traditional way.

\paragraph{Pansychism and Animism}


Recently I heard a biologist on the radio claim that biology did not have a good definition for ``life''. By that, he meant that the transition from ``inanimate matter'' to life was not clear cut. Curiously, \textbf{Rene Guenon} makes the identical point, although with a totally different meaning. In the chapter on Shamanism in the \emph{Reign of Quantity}, he writes:


\begin{quotex}
there can in fact be no `inanimate' objects in existence, and also that `life' is one of the conditions to which all corporeal existence without exception is subject; and that is why nobody has ever arrived at a satisfactory definition of the difference between the `living' and the `non-living', for that question, like so many others in modern philosophy and science, is only insoluble because there is no good reason for posing it, since the `nonliving' has no place in the domain to which the question is related, and the only differences involved are really no more than mere differences of degree.
\end{quotex}


By that, Guenon opposes the ``animism'' of the shamans to the ``mechanism'' of profane biologists. Thus, there is a ``psychic'' or ``animic'' aspect to corporeal manifestation. Another way to phrase it is to say that the world has an ``inside'' as well as an ``outside''. Or again, the psychic realm is related to ``quality'' as the physical realm is related to quantity.

For those philosophically inclined, \textbf{Timothy Sprigge} defends ``panpsychism'' in \emph{The Vindication of Absolute Idealism}, although from a different perspective.

\paragraph{Spirit in the World}


In \textbf{Karl Rahner}'s phrase, man is a ``spirit in the world''. Man's mode of being in the world, then, is non-dual: not a soul trapped in matter, or a mind interacting with inanimate matter. Rather, man participates in the world directly. Matter is what individuates one person from another quantitatively, but the experience of quality (``qualia'') is direct, since the material world is itself ensouled.

Hence, our experience of gravity is not just quantitative as in physics, but we also experience the pull of gravity directly. A thermometer will show the temperature quantitatively, but we have a direct inner experience of hot and cold.

An Indian colleague recently asked me if I could experience what a plant experiences or a cat. Not exactly, I answered, since I cannot experience the subjectivity of a \emph{particular} plant or cat. Nevertheless, I can experience what it is like to be a plant or an animal. Man incorporates and integrates the three souls of the material world. Hence, we can experience growth, respiration, the attraction to light, and so on, like a plant. Continuing, we also experience hunger, desire, fear, satiation, in short, anything an animal might experience.

This is all part of corporeal life: not merely subjective, but an authentic way of being in the world.

\paragraph{God's Knowledge of the World}


Since God is not corporeal, he does not have the experience of the world that man does. Qualia like color, warmth, hunger, and so on are all corporeal as mentioned above, that is, in the broad sense of the corporeal comprising both the physical and psychical elements. A fortiori, as the Eye of Providence, he is not looking out over the world in an exterior sense. Rather, he is at the absolute center of the world.

In a sense, this is what \textbf{Carl Jung} was getting at in the \emph{Answer to Job} that God becomes conscious through man.

Specifically, through man, he is experiencing the qualia of creation. This may give some insight into the meaning and purpose of the incarnation.

So God knows the essences of things and through this knowing, they are brought into existence. This is the identity of knowing and being. Another way to put it is that he knows things because he wills them, analogous to the way that you know what you're having for lunch at the deli: you will your choice, and don't need a scientific theory to predict what you'll get for lunch.

God's knowledge is ``all at once'' and does not need time. By knowing the possibilities of manifestation in his mind, he thereby brings them into manifestation. The world is a great unfolding, or ``evolution'' in the literal sense, of these possibilities.


\flrightit{Posted on 2015-07-29 by Cologero}
